\clearpage\thispagestyle{otFront}\phantomsection\label{mode-emploi}\pdfbookmark[0]{Mode d’emploi}{bm-mode-emploi}\otPageTitre{Mode d’emploi}{Des repères pour travailler seul ou accompagné}\fontsize{11}{15}\selectfont

\begin{tabularx}{\linewidth}{@{}lY@{}}\toprule
Une pastille & Découverte : appliquer directement le cours.\\[3mm]
Deux pastilles & Entraînement : choisir une méthode, organiser ses étapes.\\[3mm]
Trois pastilles & Problèmes : comprendre, modéliser, calculer, répondre.\\\bottomrule
\end{tabularx}\par\medskip
Ces pastilles décrivent les tâches ; elles ne classent pas les élèves.\par\medskip
\begin{exemple}[Lire une correction]Pour « $7\times\square=3$ », on utilise la
définition du quotient : le nombre caché est $\frac37$. Contrôle :
$7\times\frac37=3$. Refaire ensuite le raisonnement sans regarder aide
plus que recopier la réponse.\end{exemple}
\begin{attention}Lire l’énoncé en entier avant de calculer. Un mot comme
« plus » ou « moins » ne décide pas de l’opération : c’est la situation
qui la choisit. En géométrie, les consignes nomment toujours les objets :
la droite (AB), le segment [AB], la demi-droite [AB).\end{attention}
Chercher d’abord, ouvrir le corrigé ensuite. Après lecture, refaire le
passage difficile sur une feuille : c’est là que l’on apprend.

\clearpage
